\chapter{Bailey transforms and the Andrews--Gordon identities}\label{chap:bailey-andrews-gordon}

\index{Bailey lemma}
\index{Bailey chain}
\index{Andrews--Gordon identities}
\index{Rogers--Ramanujan continued fraction}

The limiting Bailey lemma and the two Rogers--Ramanujan series--product
identities already appear, with proofs, in
\Cref{chap:bailey}.  This chapter does not duplicate those results.  It adds
the genuinely stronger nonlimiting Bailey transform, the parameter-lowering
operation which moves between Bailey chains, the combinatorial and
continued-fraction forms of Rogers--Ramanujan theory, and every member of the
Andrews--Gordon family.

Throughout this chapter,
\[
  q\in\ComplexNumbers,\qquad 0<\AbsoluteValue q<1,
\]
and \((a_1,\ldots,a_s;q)_n\) abbreviates
\(\prod_{\nu=1}^{s}(a_\nu;q)_n\).  Thus \((q;q)_n\ne0\).
As in \Cref{chap:bailey}, the basic relation is the following.

\begin{definition}[Bailey pair]
\label{qg:def-bailey-pair}
Sequences \((\alpha_n)_{n\ge0}\) and
\((\beta_n)_{n\ge0}\) form a \emph{Bailey pair relative to \(a\)} when
\begin{equation}
  \beta_n
  =\sum_{r=0}^{n}
    \frac{\alpha_r}{(q;q)_{n-r}(aq;q)_{n+r}}
  \qquad(n\ge0).
  \label{qg:eq:bailey-pair-recall}
\end{equation}
\end{definition}

Whenever a quotient is displayed below, its denominator is assumed nonzero.
The individual statements spell out the only parameter restrictions beyond
\(0<|q|<1\).

\section{The full Bailey lemma}
\label{qg:sec:full-bailey-lemma}

\begin{theorem}[Full Bailey lemma]
\label{qg:thm-bailey-lemma}
Let \(a,\rho_1,\rho_2\in\ComplexNumbers\), with
\(\rho_1\rho_2\ne0\), and suppose that
\[
  (aq;q)_m,\qquad
  (aq/\rho_1;q)_m,\qquad
  (aq/\rho_2;q)_m
\]
are nonzero for every \(m\ge1\).  Let \((\alpha,\beta)\) be a
Bailey pair relative to \(a\), and put
\[
  A=\frac{aq}{\rho_1\rho_2}.
\]
Define
\begin{align}
 \alpha_n^\star
 &=
 \frac{(\rho_1,\rho_2;q)_n A^n}
      {(aq/\rho_1,aq/\rho_2;q)_n}\,\alpha_n,
 \label{qg:eq:full-bailey-alpha}\\
 \beta_n^\star
 &=
 \sum_{j=0}^{n}
 \frac{(\rho_1,\rho_2;q)_j(A;q)_{n-j}A^j}
      {(q;q)_{n-j}(aq/\rho_1,aq/\rho_2;q)_n}\,\beta_j.
 \label{qg:eq:full-bailey-beta}
\end{align}
Then \((\alpha^\star,\beta^\star)\) is a Bailey pair relative to
\(a\).
\end{theorem}

\begin{proof}
Insert \eqref{qg:eq:bailey-pair-recall} into
\eqref{qg:eq:full-bailey-beta} and interchange the two finite sums.
For fixed \(r\), set
\[
  N=n-r,\qquad j=r+k,\qquad c=aq^{2r+1}.
\]
The coefficient of \(\alpha_r\) is
\begin{align}
 &\frac{(\rho_1,\rho_2;q)_rA^r}
 {(aq/\rho_1,aq/\rho_2;q)_n(aq;q)_{2r}}
 \sum_{k=0}^{N}
 \frac{(\rho_1q^r,\rho_2q^r;q)_k(A;q)_{N-k}A^k}
 {(q;q)_k(q;q)_{N-k}(c;q)_k}.
 \label{qg:eq:full-bailey-inner}
\end{align}
The finite reversal identity
\begin{equation}
 \frac{(A;q)_{N-k}}{(q;q)_{N-k}}
 =
 \frac{(A;q)_N}{(q;q)_N}
 \frac{(q^{-N};q)_k}{(q^{1-N}/A;q)_k}
 \left(\frac qA\right)^k
 \label{qg:eq:finite-reversal}
\end{equation}
follows by reversing the last \(k\) factors in the two products.
Consequently the sum in \eqref{qg:eq:full-bailey-inner} equals
\[
 \frac{(A;q)_N}{(q;q)_N}
 {}_3\phi_2\!\left[
 \begin{matrix}
   \rho_1q^r,\rho_2q^r,q^{-N}\\
   c,q^{1-N}/A
 \end{matrix};q,q
 \right].
\]
This series is balanced: if \(x=\rho_1q^r\) and
\(y=\rho_2q^r\), then
\[
  \frac{xyq^{1-N}}c
  =\frac{\rho_1\rho_2q^{-N}}a
  =\frac{q^{1-N}}A.
\]
The already proved \(q\)-Pfaff--Saalsch\"utz summation
\Cref{thm:q-pfaff} therefore evaluates the last display as
\begin{equation}
  \frac{(aq^{r+1}/\rho_1,aq^{r+1}/\rho_2;q)_N}
       {(q;q)_N(c;q)_N}.
  \label{qg:eq:full-bailey-inner-evaluation}
\end{equation}

Now use the product splittings
\[
 (aq/\rho_\nu;q)_n
 =(aq/\rho_\nu;q)_r(aq^{r+1}/\rho_\nu;q)_N
 \quad(\nu=1,2)
\]
and
\[
  (aq;q)_{2r}(aq^{2r+1};q)_N=(aq;q)_{n+r}.
\]
After cancellation, the coefficient
\eqref{qg:eq:full-bailey-inner} becomes
\[
 \frac{1}{(q;q)_{n-r}(aq;q)_{n+r}}
 \frac{(\rho_1,\rho_2;q)_rA^r}
      {(aq/\rho_1,aq/\rho_2;q)_r}.
\]
The second factor is exactly the multiplier which defines
\(\alpha_r^\star\) from \(\alpha_r\) in
\eqref{qg:eq:full-bailey-alpha}.  Thus, after multiplication by
\(\alpha_r\), the coefficient contribution is the Bailey kernel times
\(\alpha_r^\star\).  Summing over \(r\) gives
\[
  \beta_n^\star
  =\sum_{r=0}^{n}
    \frac{\alpha_r^\star}
         {(q;q)_{n-r}(aq;q)_{n+r}},
\]
which is the required Bailey-pair relation.

The reversal and Pfaff--Saalsch\"utz calculation was written on the
Zariski-open set where all temporary denominators are nonzero.  On that set
it proves an identity of rational functions in \(a,\rho_1,\rho_2\).
Multiplying by a common denominator turns it into a polynomial identity, so
it holds identically.  Therefore the coefficient formula extends to every
parameter choice satisfying the permanent nonvanishing assumptions in the
theorem, including choices at which one of the temporary factors vanishes.
\end{proof}

\begin{remark}[The limiting theorem is not restated]
\label{qg:rem:bailey-limiting-reference}
Letting \(\rho_1,\rho_2\) tend to infinity in
\Cref{qg:thm-bailey-lemma}, with \(n\) fixed, gives
\[
 \frac{(\rho_1,\rho_2;q)_n}
      {(aq/\rho_1,aq/\rho_2;q)_n}
 \left(\frac{aq}{\rho_1\rho_2}\right)^n
 \longrightarrow a^nq^{n^2}.
\]
The beta summands have the corresponding finite limits.  The resulting pair
transformation is precisely \Cref{thm:bailey-limit-finite}; its absolutely
convergent infinite consequence is already
\Cref{cor:bailey-limit-infinite}.  Thus the present theorem strictly supplies
the two free parameters which the backbone's limiting form omits.
\end{remark}

\section{Lowering the relative parameter}
\label{qg:sec:bailey-parameter-lowering}

\begin{lemma}[Parameter-lowering shift]
\label{qg:lem-bailey-lowering}
Let \((\alpha,\beta)\) be a Bailey pair relative to \(a\), and assume
\((a;q)_m\ne0\) for every \(m\ge1\).  Define
\begin{align}
 \widetilde\alpha_0&=\alpha_0,
 \label{qg:eq:bailey-lowering-zero}\\
 \widetilde\alpha_n
 &=(1-a)\left(
 \frac{\alpha_n}{1-aq^{2n}}
 -\frac{aq^{2n-2}\alpha_{n-1}}{1-aq^{2n-2}}
 \right)
 \qquad(n\ge1),
 \label{qg:eq:bailey-lowering-positive}\\
 \widetilde\beta_n&=\beta_n.
 \label{qg:eq:bailey-lowering-beta}
\end{align}
Then \((\widetilde\alpha,\widetilde\beta)\) is a Bailey pair
relative to \(a/q\).
\end{lemma}

\begin{proof}
A Bailey pair relative to \(a/q\) has kernel
\[
  \frac{1}{(q;q)_{n-j}(a;q)_{n+j}}.
\]
Substitute \eqref{qg:eq:bailey-lowering-zero} and
\eqref{qg:eq:bailey-lowering-positive} into the sum against this
kernel, and reindex the part containing \(\alpha_{j-1}\).  For fixed
\(0\le r\le n\), the coefficient of \(\alpha_r\) is
\begin{align*}
 &\frac{1-a}{1-aq^{2r}}\,
  \frac{1}{(q;q)_{n-r}(a;q)_{n+r}}\\
 &\quad-
 \frac{(1-a)aq^{2r}}{1-aq^{2r}}\,
 \frac{1}{(q;q)_{n-r-1}(a;q)_{n+r+1}},
\end{align*}
where the second term is absent if \(r=n\).  Put both terms over
\[
 (q;q)_{n-r}(a;q)_{n+r+1}(1-aq^{2r}).
\]
Their numerator, apart from the common factor \(1-a\), is
\[
 (1-aq^{n+r})-aq^{2r}(1-q^{n-r})
 =1-aq^{2r}.
\]
This formula also covers \(r=n\), because then \(1-q^{n-r}=0\).
After cancellation, the coefficient is
\[
 \frac{1-a}{(q;q)_{n-r}(a;q)_{n+r+1}}
 =
 \frac{1}{(q;q)_{n-r}(aq;q)_{n+r}}.
\]
It is exactly the original Bailey kernel.  Therefore the transformed sum is
\(\beta_n=\widetilde\beta_n\), proving the claim.
\end{proof}

\begin{remark}
\Cref{qg:lem-bailey-lowering} is more general than the
specialization needed below.  With \(a=q\), it changes a pair relative to
\(q\) into a pair relative to \(1\) without changing its beta sequence.
Moving that single shift to different depths of a Bailey chain is what
creates the different linear tails in the Andrews--Gordon sums.
\end{remark}

\section{Iterating the limiting chain}
\label{qg:sec:iterated-bailey-chain}

Let \(S_a\) denote exactly the finite pair transformation proved in
\Cref{thm:bailey-limit-finite}.  We use a superscript \(^{(t)}\) for
the result of applying \(S_a\) \(t\) times.

\begin{lemma}[Iterated Bailey chain]
\label{qg:lem-bailey-chain-iterate}
For every Bailey pair \((\alpha,\beta)\) relative to \(a\) and every
\(t\in\NaturalNumbers\), put
\((\alpha^{(0)},\beta^{(0)})=(\alpha,\beta)\).  For \(t\ge1\),
\begin{align}
 \alpha_n^{(t)}
 &=a^{tn}q^{tn^2}\alpha_n,
 \label{qg:eq:iterated-bailey-alpha}\\
 \beta_n^{(t)}
 &=
 \sum_{n\ge r_1\ge\cdots\ge r_t\ge0}
 \frac{
   a^{r_1+\cdots+r_t}
   q^{r_1^2+\cdots+r_t^2}\beta_{r_t}}
 {(q;q)_{n-r_1}
  \displaystyle\prod_{\ell=1}^{t-1}
  (q;q)_{r_\ell-r_{\ell+1}}}.
 \label{qg:eq:iterated-bailey-beta}
\end{align}
In particular, if \(\beta_n=\delta_{n,0}\), then
\(\beta_n^{(1)}=1/(q;q)_n\), while for \(t\ge2\),
\begin{equation}
 \beta_n^{(t)}
 =
 \sum_{n\ge r_1\ge\cdots\ge r_{t-1}\ge0}
 \frac{
   a^{r_1+\cdots+r_{t-1}}
   q^{r_1^2+\cdots+r_{t-1}^2}}
 {(q;q)_{n-r_1}
  \displaystyle\prod_{\ell=1}^{t-2}
  (q;q)_{r_\ell-r_{\ell+1}}
  (q;q)_{r_{t-1}}},
 \label{qg:eq:iterated-bailey-unit-beta}
\end{equation}
\end{lemma}

\begin{proof}
The case \(t=0\) is tautological.  One application of
\Cref{thm:bailey-limit-finite} multiplies the alpha component by
\(a^nq^{n^2}\), which proves
\eqref{qg:eq:iterated-bailey-alpha} by induction.

For the beta component, one application of the same transformation gives
\eqref{qg:eq:iterated-bailey-beta} at depth \(1\), where the product is
empty.  Suppose the formula holds at some depth \(t\ge1\).  The next
application of \(S_a\) gives
\[
 \beta_n^{(t+1)}
 =\sum_{j=0}^{n}
   \frac{a^jq^{j^2}}{(q;q)_{n-j}}\,
   \beta_j^{(t)}.
\]
Insert the induction hypothesis with \(j\) as its outer bound and rename
\(j=r_1\).  This adds precisely the new outer inequality
\(n\ge r_1\), the factor
\(a^{r_1}q^{r_1^2}\), and the denominator
\((q;q)_{n-r_1}\), yielding
\eqref{qg:eq:iterated-bailey-beta} at depth \(t+1\).
Finally, if \(\beta_{r_t}=\delta_{r_t,0}\), only \(r_t=0\)
survives; deleting that forced index gives
\eqref{qg:eq:iterated-bailey-unit-beta}.
\end{proof}

\section{Rogers--Ramanujan partitions and the continued fraction}
\label{qg:sec:rogers-ramanujan-forms}

Write
\[
 G(q)=\sum_{n\ge0}\frac{q^{n^2}}{(q;q)_n},
 \qquad
 H(q)=\sum_{n\ge0}\frac{q^{n^2+n}}{(q;q)_n}.
\]
Their product evaluations are exactly
\Cref{thm:rogers-ramanujan}, so we use rather than restate that theorem.

\begin{corollary}[Rogers--Ramanujan partition interpretations]
\label{qg:cor-rr-partitions}
For every \(N\ge0\):
\begin{enumerate}[label=\textup{(\roman*)}]
\item the number of partitions of \(N\) whose successive parts differ by
at least \(2\) equals the number of partitions of \(N\) into parts
congruent to \(1\) or \(4\pmod 5\);
\item the number of partitions of \(N\) whose successive parts differ by
at least \(2\) and whose smallest part is at least \(2\) equals the number
of partitions of \(N\) into parts congruent to \(2\) or \(3\pmod 5\).
\end{enumerate}
\end{corollary}

\begin{proof}
Fix the number \(n\) of parts.  A partition with successive differences at
least \(2\) contains the staircase
\[
  (2n-1,2n-3,\ldots,3,1),
\]
whose size is \(n^2\).  Subtracting it componentwise leaves an arbitrary
weakly decreasing \(n\)-tuple of nonnegative integers, equivalently a
partition with at most \(n\) parts.  By
\Cref{prop:bounded-partitions}, its generating function is
\(1/(q;q)_n\).  Thus the generating function for the partitions in
part~\textup{(i)} is \(G(q)\).

If the smallest part must be at least \(2\), subtract instead
\[
  (2n,2n-2,\ldots,4,2),
\]
whose size is \(n^2+n\).  The remainder is again an arbitrary partition
with at most \(n\) parts, so the generating function is \(H(q)\).

By \Cref{thm:rogers-ramanujan}, \(G(q)\) is the reciprocal product over
positive integers congruent to \(1\) or \(4\pmod5\), and \(H(q)\) is
the reciprocal product over those congruent to \(2\) or \(3\pmod5\).
Each reciprocal factor permits an arbitrary multiplicity of its part.
All these products and series are holomorphic at \(q=0\) throughout the
open unit disk.  The analytic identities there therefore imply, by
uniqueness of Taylor expansions, equality coefficient by coefficient.
This proves both claims.
\end{proof}

For each \(m\in\NaturalNumbers\), introduce the single family
\begin{equation}
  F_m(q)=\sum_{n\ge0}\frac{q^{n^2+mn}}{(q;q)_n}.
  \label{qg:eq:rogers-family}
\end{equation}
Thus \(F_0=G\) and \(F_1=H\).

\begin{lemma}[Rogers recurrence]
\label{qg:lem-rogers-recurrence}
For every \(m\in\NaturalNumbers\),
\begin{equation}
  F_m(q)=F_{m+1}(q)+q^{m+1}F_{m+2}(q).
  \label{qg:eq:rogers-recurrence}
\end{equation}
\end{lemma}

\begin{proof}
The \(n=0\) terms cancel when \(F_m-F_{m+1}\) is formed.  For
\(n\ge1\), use
\((1-q^n)/(q;q)_n=1/(q;q)_{n-1}\):
\begin{align*}
 F_m-F_{m+1}
 &=\sum_{n\ge1}
   \frac{q^{n^2+mn}(1-q^n)}{(q;q)_n}\\
 &=\sum_{n\ge1}\frac{q^{n^2+mn}}{(q;q)_{n-1}}\\
 &=q^{m+1}\sum_{r\ge0}
   \frac{q^{r^2+(m+2)r}}{(q;q)_r}
 =q^{m+1}F_{m+2},
\end{align*}
where \(r=n-1\).  Absolute convergence justifies the subtraction.
\end{proof}

\begin{theorem}[Rogers--Ramanujan continued fraction]
\label{qg:thm-rrcf}
For \(|q|<1\), the single-valued ratio satisfies
\begin{equation}
 \frac{F_1(q)}{F_0(q)}
 =
 \cfrac{1}{
  1+\cfrac{q}{
  1+\cfrac{q^2}{
  1+\cfrac{q^3}{1+\ddots}}}},
 \label{qg:eq:rogers-ramanujan-fraction}
\end{equation}
and the continued fraction converges locally uniformly in the unit disk.
On any simply connected subdomain of \(0<|q|<1\), choose a branch of
\(q^{1/5}\) and define
\begin{equation}
 R(q)
 \coloneqq q^{1/5}\frac{F_1(q)}{F_0(q)}
 =
 q^{1/5}
 \frac{(q,q^4;q^5)_\infty}
      {(q^2,q^3;q^5)_\infty}.
 \label{qg:eq:rogers-ramanujan-fraction-product}
\end{equation}
\end{theorem}

\begin{proof}
For \(N\in\NaturalNumbers\), consider the finite continuant
\[
 C_N
 =1+\cfrac{q}{
  1+\cfrac{q^2}{
  \ddots+\cfrac{q^N}{1}}}
 =\frac{P_N}{Q_N}.
\]
Its numerator and denominator satisfy
\begin{align*}
 P_{-1}=1,\quad P_0=1,\quad
 P_N&=P_{N-1}+q^NP_{N-2},\\
 Q_{-1}=0,\quad Q_0=1,\quad
 Q_N&=Q_{N-1}+q^NQ_{N-2}.
\end{align*}
We claim
\begin{align}
 P_N
 &=\sum_{j=0}^{\lfloor(N+1)/2\rfloor}
   q^{j^2}\GaussianBinomial{N+1-j}{j}{q},
 \label{qg:eq:continuant-P}\\
 Q_N
 &=\sum_{j=0}^{\lfloor N/2\rfloor}
   q^{j^2+j}\GaussianBinomial{N-j}{j}{q}.
 \label{qg:eq:continuant-Q}
\end{align}
The formulas give the displayed initial values.  For
\eqref{qg:eq:continuant-P}, apply the second \(q\)-Pascal recurrence
from \Cref{thm:qpascal}:
\[
 \GaussianBinomial{N+1-j}{j}{q}
 =
 \GaussianBinomial{N-j}{j}{q}
 +q^{N+1-2j}\GaussianBinomial{N-j}{j-1}{q}.
\]
After multiplication by \(q^{j^2}\), the first term sums to
\(P_{N-1}\).  In the second term put \(h=j-1\); its power is
\[
 j^2+N+1-2j=N+h^2,
\]
so it sums to \(q^NP_{N-2}\).  This proves
\eqref{qg:eq:continuant-P} by induction.  The same argument, using
\[
 \GaussianBinomial{N-j}{j}{q}
 =
 \GaussianBinomial{N-1-j}{j}{q}
 +q^{N-2j}\GaussianBinomial{N-1-j}{j-1}{q},
\]
proves \eqref{qg:eq:continuant-Q}; after \(h=j-1\), the second
exponent is \(N+h^2+h\).

For fixed \(j\),
\[
 \GaussianBinomial{N+1-j}{j}{q}
 =\frac{(q^{N+2-2j};q)_j}{(q;q)_j}
 \longrightarrow\frac1{(q;q)_j},
\]
and the analogous limit holds in
\eqref{qg:eq:continuant-Q}.  The convergence may be passed through
the sums locally uniformly.  Indeed, on \(|q|\le r<1\), every Gaussian
coefficient occurring above is bounded in modulus by
\[
 C_r=\frac{(-r;r)_\infty}{(r;r)_\infty},
\]
because all powers in its numerator are positive; the remaining factors
\(r^{j^2}\) and \(r^{j^2+j}\) are summable.  Dominated convergence gives
\[
  P_N\longrightarrow F_0(q),
  \qquad
  Q_N\longrightarrow F_1(q)
\]
locally uniformly.

The products in \Cref{thm:rogers-ramanujan} show that \(F_0\) and
\(F_1\) are nonzero for \(|q|<1\).  Let \(K_0\) be a compact subset of
the unit disk.  Continuity and zero-freeness give
\[
  \eta_{K_0}=\min_{q\in K_0}|F_0(q)|>0.
\]
Uniform convergence \(P_N\to F_0\) gives
\(|P_N-F_0|<\eta_{K_0}/2\) on \(K_0\) for all sufficiently large
\(N\).  Those \(P_N\) are consequently zero-free on \(K_0\), and
\[
  C_N^{-1}=\frac{Q_N}{P_N}
  \longrightarrow\frac{F_1(q)}{F_0(q)}.
\]
The convergence is uniform on \(K_0\).  We define each finite truncation by
the rational continuant \(Q_N/P_N\); where the nested expression is defined,
the two agree, and the continuant supplies its analytic continuation across
removable intermediate singularities.  This proves
\eqref{qg:eq:rogers-ramanujan-fraction} and its local uniform convergence.
Finally, on any chosen fifth-root branch, taking the ratio of the two already
proved product evaluations in \Cref{thm:rogers-ramanujan} and multiplying by
\(q^{1/5}\) gives
\eqref{qg:eq:rogers-ramanujan-fraction-product}.
\end{proof}

\begin{remark}[Schur-type finitizations]
The polynomials \(P_N\) and \(Q_N\) in
\eqref{qg:eq:continuant-P}--\eqref{qg:eq:continuant-Q} are Schur-type
finitizations of \(F_0\) and \(F_1\).  Thus the same recurrence has three
compatible realizations: a finite rational continuant, an exact finite sum of
Gaussian polynomials, and, after \(N\to\infty\), a Rogers--Ramanujan series.
The finite level is useful in its own right: it retains coefficient data,
exhibits positivity, and permits exact induction before any analytic limit is
taken.  This is a reusable proof pattern.  When a quotient of \(q\)-series is
governed by a short recurrence, one should first look for its continuants and
then for Gaussian-polynomial formulas for those continuants; convergence and
the infinite identity may then emerge as the final step of finite algebra.
\end{remark}

\section{The full Andrews--Gordon family}
\label{qg:sec:andrews-gordon}

\begin{theorem}[Andrews--Gordon identities]
\label{qg:thm-andrews-gordon}
Let \(k,i\in\NaturalNumbers\) satisfy \(k\ge2\) and \(1\le i\le k\).  For
\(n_1,\ldots,n_{k-1}\ge0\), put
\[
  N_j=n_j+n_{j+1}+\cdots+n_{k-1}
  \qquad(1\le j\le k-1).
\]
Then
\begin{equation}
 \sum_{n_1,\ldots,n_{k-1}\ge0}
 \frac{
 q^{N_1^2+\cdots+N_{k-1}^2+N_i+\cdots+N_{k-1}}}
 {(q;q)_{n_1}\cdots(q;q)_{n_{k-1}}}
 =
 \frac{(q^i,q^{2k+1-i},q^{2k+1};q^{2k+1})_\infty}
      {(q;q)_\infty},
 \label{qg:eq:andrews-gordon}
\end{equation}
where \(N_i+\cdots+N_{k-1}\) is the empty sum \(0\) when \(i=k\).
\end{theorem}

\begin{proof}
Put
\[
  K=k-1,\qquad M=2k+1,
\]
and use decreasing variables
\[
  r_j=N_j\quad(1\le j\le K),\qquad r_{K+1}=0.
\]
Then \(n_j=r_j-r_{j+1}\), so the left side of
\eqref{qg:eq:andrews-gordon} is
\begin{equation}
 \sum_{r_1\ge\cdots\ge r_K\ge0}
 \frac{
 q^{r_1^2+\cdots+r_K^2+r_i+\cdots+r_K}}
 {(q;q)_{r_1-r_2}\cdots
  (q;q)_{r_{K-1}-r_K}(q;q)_{r_K}}.
 \label{qg:eq:andrews-gordon-decreasing}
\end{equation}
We derive this expression from the unit pair relative to \(q\) in
\Cref{cor:unit-bailey-pairs}.

\smallskip
\noindent\emph{The case \(i=1\).}
Apply \(S_q\) exactly \(K\) times, and then apply the infinite limiting
identity \Cref{cor:bailey-limit-infinite}, still relative to \(q\).
By \Cref{qg:lem-bailey-chain-iterate}, the beta side has \(K\)
decreasing variables after the final outer summation.  Every one has weight
\(q^{r_j^2+r_j}\); its denominator is exactly that in
\eqref{qg:eq:andrews-gordon-decreasing}.  Hence the beta side is the
desired sum for \(i=1\).

The unit alpha sequence relative to \(q\), from
\Cref{cor:unit-bailey-pairs}, is
\[
 \alpha_n^{\mathrm U}
 =(-1)^nq^{\binom n2}\frac{1-q^{2n+1}}{1-q}.
\]
The \(K\) chain steps and the final limiting identity together multiply
this by \(q^{k(n^2+n)}\).  Since
\((q^2;q)_\infty=(q;q)_\infty/(1-q)\), the alpha side is
\begin{align*}
 \frac1{(q;q)_\infty}
 \sum_{n\ge0}(-1)^n
 q^{\{Mn^2+(M-2)n\}/2}(1-q^{2n+1})
 &=
 \frac1{(q;q)_\infty}
 \sum_{m\in\IntegerNumbers}(-1)^m
 q^{\{Mm^2+(M-2)m\}/2}.
\end{align*}
The term containing \(-q^{2n+1}\) becomes the negative-index term
\(m=-n-1\).  Apply the Jacobi triple product
\Cref{thm:jtp} with base \(q^M\) and parameter \(q^{M-1}\).
The bilateral sum is
\[
  (q^{M-1},q,q^M;q^M)_\infty,
\]
which is the numerator in \eqref{qg:eq:andrews-gordon} for
\(i=1\).

\smallskip
\noindent\emph{The cases \(2\le i\le k\).}
Set
\[
  t=k-i+1,\qquad u=i-2.
\]
Thus \(t\ge1\), \(u\ge0\), and \(t+u=K\).  Starting with the same
unit pair relative to \(q\), perform, in order,
\[
  \underbrace{S_q\circ\cdots\circ S_q}_{t\ \mathrm{times}},
  \qquad
  \text{the lowering shift of
        \Cref{qg:lem-bailey-lowering}},
  \qquad
  \underbrace{S_1\circ\cdots\circ S_1}_{u\ \mathrm{times}},
\]
and finally apply \Cref{cor:bailey-limit-infinite} relative to \(1\).
The lowering shift leaves beta unchanged.  The first \(S_q\), applied to
the delta beta sequence, creates the terminal factor \(1/(q;q)_n\);
the remaining \(t-1=k-i\) \(q\)-steps create the innermost variables,
with weights \(q^{r_j^2+r_j}\).  The \(u\) subsequent \(1\)-steps and
the final limiting sum create the outer \(u+1=i-1\) variables, with
weights \(q^{r_j^2}\).  Thus the linear weights occur exactly for
\(r_i,\ldots,r_K\), and the beta side is
\eqref{qg:eq:andrews-gordon-decreasing}.

It remains to compute the alpha side.  After the first \(t\) steps, write
\[
 \gamma_n=q^{t(n^2+n)}\alpha_n^{\mathrm U}.
\]
The lowering formula gives \(\delta_0=1\), while for \(n\ge1\),
\begin{align}
 \delta_n
 &=(1-q)\left(
   \frac{\gamma_n}{1-q^{2n+1}}
   -\frac{q^{2n-1}\gamma_{n-1}}{1-q^{2n-1}}
   \right)\notag\\
 &=(-1)^n
 q^{\{(2t+1)n^2-(2t-1)n\}/2}
 \left(1+q^{(2t-1)n}\right).
 \label{qg:eq:andrews-gordon-lowered-alpha}
\end{align}
For completeness, the two powers before factoring in the last line are
\[
 t(n^2+n)+\binom n2
 \quad\text{and}\quad
 2n-1+t(n^2-n)+\binom{n-1}{2};
\]
their difference is \((2t-1)n\), which proves
\eqref{qg:eq:andrews-gordon-lowered-alpha}.

The \(u\) transformations \(S_1\) and the final limiting identity
multiply \(\delta_n\) by \(q^{(u+1)n^2}=q^{(i-1)n^2}\).
Since
\[
  2(t+i)-1=M,\qquad 2t-1=M-2i,
\]
the alpha side becomes
\begin{align*}
 \frac1{(q;q)_\infty}
 \left(
 1+\sum_{n\ge1}(-1)^n
 q^{\{Mn^2-(M-2i)n\}/2}
 \left(1+q^{(M-2i)n}\right)
 \right)
 &=
 \frac1{(q;q)_\infty}
 \sum_{m\in\IntegerNumbers}(-1)^m
 q^{\{Mm^2-(M-2i)m\}/2}.
\end{align*}
Here the two positive-index terms pair \(m=n\) with \(m=-n\).
The Jacobi triple product with base \(q^M\) and parameter \(q^i\)
evaluates the bilateral sum as
\[
  (q^i,q^{M-i},q^M;q^M)_\infty.
\]
This proves \eqref{qg:eq:andrews-gordon} in every case.

All finite Bailey transformations above are finite rearrangements.  For
\(|q|\le R<1\),
\[
  |(q;q)_m|
  \ge\prod_{\nu=1}^{m}(1-R^\nu)
  \ge(R;R)_\infty,
\]
uniformly in \(m\).  Since \(K=k-1\) is fixed, every multiple-series
summand is therefore bounded by a constant depending only on \(R\) and
\(K\), times \(R^{r_1^2+\cdots+r_K^2}\).  Dropping the monotonicity
conditions gives the summable majorant
\[
  C_{R,K}\left(\sum_{n\ge0}R^{n^2}\right)^K.
\]
The alpha and bilateral sums are bounded similarly by
\(C'_R R^{c n^2-Cn}\) with \(c>0\).  Hence all final series converge
absolutely and locally uniformly, and every limiting passage used in the
proof satisfies the convergence hypothesis of
\Cref{cor:bailey-limit-infinite}.
\end{proof}

\begin{corollary}[The modulus-seven specialization]
\label{qg:cor-andrews-gordon-mod7}
For \(r,s\ge0\), put \(N_1=r+s\) and \(N_2=s\).  Then
\begin{align}
 \sum_{r,s\ge0}
 \frac{q^{N_1^2+N_2^2+N_1+N_2}}
      {(q;q)_r(q;q)_s}
 &=
 \frac1{(q^2,q^3,q^4,q^5;q^7)_\infty},
 \label{qg:eq:andrews-gordon-mod-seven-one}\\
 \sum_{r,s\ge0}
 \frac{q^{N_1^2+N_2^2+N_2}}
      {(q;q)_r(q;q)_s}
 &=
 \frac1{(q,q^3,q^4,q^6;q^7)_\infty},
 \label{qg:eq:andrews-gordon-mod-seven-two}\\
 \sum_{r,s\ge0}
 \frac{q^{N_1^2+N_2^2}}
      {(q;q)_r(q;q)_s}
 &=
 \frac1{(q,q^2,q^5,q^6;q^7)_\infty}.
 \label{qg:eq:andrews-gordon-mod-seven-three}
\end{align}
Equivalently, the coefficient of \(q^n\) in the \(i\)-th sum
\((i=1,2,3)\) counts partitions of \(n\) into parts not congruent to
\(0\) or \(\pm i\pmod7\).
\end{corollary}

\begin{proof}
Take \(k=3\) and respectively \(i=1,2,3\) in
\Cref{qg:thm-andrews-gordon}.  The three numerator products are
\[
 (q,q^6,q^7;q^7)_\infty,\qquad
 (q^2,q^5,q^7;q^7)_\infty,\qquad
 (q^3,q^4,q^7;q^7)_\infty.
\]
Dissect
\[
 (q;q)_\infty=(q,q^2,q^3,q^4,q^5,q^6,q^7;q^7)_\infty
\]
and cancel the three displayed residue classes.  This gives
\eqref{qg:eq:andrews-gordon-mod-seven-one}--
\eqref{qg:eq:andrews-gordon-mod-seven-three}.  Each remaining reciprocal
factor \(1/(1-q^m)\) records an arbitrary multiplicity of the part \(m\),
and all these products and series are holomorphic at \(q=0\) throughout
the open unit disk.  Since the analytic identities hold there, uniqueness
of Taylor expansions proves the coefficient interpretation.
\end{proof}

\begin{remark}[Odd-modulus specialization map]
The modulus in \Cref{qg:thm-andrews-gordon} is \(2k+1\).  Consequently
\(k=4\) gives four identities of odd modulus \(9\), and \(k=13\) gives
thirteen identities of odd modulus \(27\).  These are direct specializations
of the theorem just proved and require no new product argument.  Modulus \(14\)
does not arise by ``substituting \(14\) for \(2k+1\)'': no integral \(k\)
satisfies that equation.  Its natural setting is the parity-sensitive
even-modulus theory of Bressoud, whose extra parity condition is absent from
the Andrews--Gordon sums above.  Constructing an equally elementary
Bailey-lattice route to that theory is the
\hyperref[qg:prob-even-modulus-lowering]{even-modulus lowering open problem}
recorded in the next chapter.
\end{remark}

\begin{remark}[Logical dependencies]
\label{qg:rem:bailey-andrews-gordon-dependencies}
No external identity is being used as an unproved black box in this chapter.
The full Bailey lemma reduces to the proved
\(q\)-Pfaff--Saalsch\"utz summation \Cref{thm:q-pfaff}; the
continued fraction reduces to the proved \(q\)-Pascal recurrence and
\Cref{thm:rogers-ramanujan}; and the Andrews--Gordon products reduce to
the proved Jacobi triple product \Cref{thm:jtp}.  The limiting Bailey
results invoked above are precisely the finite and infinite theorems already
proved in \Cref{chap:bailey}.
\end{remark}
